\section{A finite-range obstruction to simultaneous witness and specialization restriction}
\label{sec:finite-restriction}
Let $\mathcal C$ be the native Choice--DI calculus with fixed scalar terms
$0,S,I$, equations, binders $Q_yB$ and $D_{x,y}B$, and the rules specified in
\texttt{ipra\_named\_specialization}. Consider the simultaneous restriction
\[
\text{DI witnesses: }0\text{ or an in-scope variable},\qquad
\text{QE targets: }1\text{ or an in-scope variable}.
\tag{R}
\]
Neither arbitrary numerical substitution nor ordinary quantified logic
is silently added as a primitive rule.

\begin{theorem}
Restriction (R) strictly decreases the native theorem set. In particular,
with $C(y,p,q)=I(y,p;u.q)$ and $1=S0$, the formula
\[
D_{x,y}\bigl(C(y,1,x)=1\bigr)
\tag{N1}
\]
is a theorem of $\mathcal C$ but not of (R).
\end{theorem}
\begin{proof}
Evaluate numerical terms by their standard natural-number functions and
assign zero to every free numerical variable. Define native formula truth
recursively by ordinary equality and
\[
\begin{split}
[Q_yB]&=[B(y:=0)]\Rightarrow[B(y:=1)],\\
[D_{x,y}B]&=[B(x:=0,y:=0)]\Rightarrow[B(x:=0,y:=1)].
\end{split}
\]
These clauses specify a proof invariant, not an ordinary first-order
interpretation with unrestricted substitution. Formula substitution
respects equality of the numerical values of substituted terms, since
all binders are renamed away from those terms before substitution.
Thus equality transport and the computation schemas preserve validity.
Hypothesis use does so immediately.

An allowed DI witness has numerical value zero. Evaluating its fresh
step parameter $z$ at zero makes the successor $Sz$ have value one.
For case DI, the premise therefore supplies exactly the implication in
the displayed $D$ interpretation. For progress DI it supplies the same
implication, because its current hypothesis is evaluated at $z=0$.
The argument even applies to compound progress; the literal-equation
restriction is not needed for this invariant. The retained context is
unaffected by the fresh parameter.

In DE the fresh parameter $a$ also has value zero. Consequently the added
hypothesis $Q_yB(a,y)$ has exactly the truth value of its major $D_{x,y}B$.
Validity of the continuation proves validity of the conclusion. A
variable-target QE concludes the same truth value as its base premise;
a target-one QE is ordinary implication elimination. All primitive rules
preserve the invariant, so external induction on finite derivations
establishes soundness.

In (N1) the zero instance is $C(0,1,0)=1$, which is true, while the one
instance is $C(1,1,0)=1$, which says $0=1$ and is false. Hence (N1) is not
derivable under (R). In the unrestricted calculus, case DI with witness
$1$ proves it immediately from the computation equation
$C(Sz,1,1)=1$, ignoring the base hypothesis.
\end{proof}

\paragraph{General finite-range version.}
Fix $N\geq1$, let $W=\{0,\ldots,N\}$, allow QE only at variables and
numerals in $W$, and continue to restrict DI witnesses to zero or
variables. Evaluate free variables in $W$, keep ordinary numerical term
semantics (whose outputs may lie outside $W$), and interpret
\[
Q_yB := B(0)\Rightarrow\bigwedge_{j\in W} B(j),\qquad
D_{x,y}B := \bigvee_{i\in W}\left(B(i,0)\Rightarrow
                              \bigwedge_{j\in W}B(i,j)\right).
\]
A sequent is valid at every assignment of its free variables in $W$.
Case DI is valid because each positive $j\leq N$ is $S(j-1)$ with
$j-1\in W$. Progress DI is valid by the finite external induction from
zero to $N$. DE uses a witness in $W$ and its continuation's validity at
that assignment; QE uses a target in $W$. Equality transport is valid
because the recursively defined formula semantics respects numerical
equality even for substituted numerical values outside $W$.
The same primitive computations hold. Thus all rules preserve validity.
The unrestricted theorem
\[
D_{x,y}\bigl(C(y,N+1,x)=N+1\bigr)
\]
is false in this interpretation, since $1\in W$ but no possible witness
$i\in W$ equals $N+1$. Adding finitely many allowed numeral targets
therefore does not repair the simultaneous restriction.

\paragraph{Scope.}
This is an exact obstruction to identity-language theorem equivalence for
these simultaneous restrictions. It is not a proof that every differently
interpreted theory with restricted rules fails to be bi-interpretable with
iPRA. The special proof invariant for (N1) is stronger than the $N=1$
finite-range model and yields a smaller separating formula.

\paragraph{Compatibility with known-zero-instance Supply.}
The counterinterpretations remain sound for the rule
\[
\frac{\Gamma,H(w)\vdash A\qquad\Gamma\vdash H(0)}
     {\Gamma\vdash A},\qquad w\notin\operatorname{FV}(\Gamma,A).
\]
In the fixed-zero invariant, $H(w)$ and $H(0)$ have the same truth value.
For the finite-range invariant, extend any outer satisfying assignment
by $w=0$, which is permitted because $0\in W$. The second premise makes
$H(w)$ true at that extension; the first premise supplies $A$. Freshness
makes the end sequent unaffected by that extension. This also covers the
new lazy computation generator, whose separate soundness theorem proves
its generated equation's zero instance. No arbitrary-term naming rule
with a witness outside $W$ is included.
